As exactly right. So in fact, you know, I'm looking for a macro-fluidic system that I can flow a series of reagents across a graphene semiconductor chip that I need to have voltage control over and ideally that's coordinated with this fluidic system. And I want to have four particular reagents at least corresponding to different DNA bases. And I also want to have, you know, the ability to push washes between them and wash out the lines on the thing. And another critical aspect is I want to be able to preserve particular washes off of the system, you know, for later analysis. So not all, like many could go to waste, but some need to be something I can pick up and work with later. So is that making sense? [Music]

So DNA synthesis is the safest way to term it.

I mean, so it should be kind of up front that the objective is DNA synthesis. But, you know, and really like looking at a new synthesis system, I won't say anything about the DNA data storage. Like, yeah.

Okay, I want to explore a new synthesis pathway. Yeah, okay.

I also want to include the spec sheet and kind of dimensions of the target liquid gated graphing field effect chips that I want to work with because that gives a sense of scale that it's not some super tiny thing. It's like I want to make a flow that goes across the whole surface of it and how big is that. [Music]

So this works then in effect there's some upside for the whole organization, which is that I could compact with a much bigger contract and a much bigger project with bigger impact. And that would be specifically, I can outline what that might look like. I can outline what that might look like. It's this, okay, you want to do individual control inside the flow, so we have to make a new chip that doesn't exist. And we have to do integration with another semiconductor system that doesn't exist, which you guys could actually make. So it kind of sets up this sort of collaborative process in a way that would be able to do it. [Music]

chatgpt-4 plus major edits

Subject: Exploring DNA Synthesis with Graphene Field Effect Transistors

Dear Nicholas,

Thanks for our discussion yesterday. I am writing to provide details related to a project opportunity exploring a new DNA synthesis modality. Practically, we would like to integrate a fluidic delivery system with liquid-gated graphene field effect transistor (gFET) chips. The objective would be to flow a series of reagents across the graphene semiconductor chip, with coordinated voltage control of the gFET channels and liquid gate, so as to ascertain the possibility of using a particular approach for DNA synthesis.

The system should accommodate at least four reagents corresponding to different DNA bases, as well as the capability to perform washes between them and clear the lines as needed, and two to three additional reagents (e.g. endonucleases) required for a full testing protocol. To ensure there is some spare capacity in the system for changes in experimental design, I suppose this would require 8-10 reservoirs and fluid input lines. Additionally, we would like to be able to preserve specific washes from the system for later biochemical analysis.

Attached to this email, you will find a data sheet of the target liquid-gated gFET chips we plan to work with, which should give a sense of the physical scale of this project. For these tests, we would like to expose a central region around 1-2mm wide to liquid, so that all the 12 gFETs are in the same environment. The manufacturer already provides a cartridge system designed for liquid sensing (data sheet attached). This could either be used as a component in this project or its general design might suggest how to interface with the gFET chips within a fluidic control platform. As an alternative to a fluid delivery system, we are contemplating using an Opentrons robot, but as we discussed, this has some disadvantages for our study.

If GTRI would be able to develop the development of the fluid delivery system and device integration with a source-measure unit required to drive and obtain signal information from the gFET, that would be very advantageous. Testing of the system through the execution of a basic synthesis protocol would also be helpful, if it could be done within reasonable budget constraints. However, I am also experienced in device integration of this type and would be able to complete the research at UTHSC. As we discussed, I'm aware of typical costs for projects of this scale and have sufficient resources to fund a contracted project with GTRI. I would also be open to collaboration with GTRI research groups, should that achieve a good balance of costs and other factors.

Successfully developing this new synthesis pathway could lead to a more extensive collaboration and a larger contract, involving the creation of a custom gFET chip (designed by myself and Graphenea) and its integration with other semiconductor systems, something which I understand GTRI has specialization doing.

Please let me know if you require any additional information or clarification on this project. I look forward to discussing this opportunity further with you and your colleagues.

Best regards,

Erik


